\pagestyle{empty}
\tight
\begin{tblr}{width=\textwidth,colspec={|Q[c,m,1.9cm]|X[l,m]|},hlines,vlines,hspan=minimal,row{1}={bg=headblue},rowsep=1pt,colsep=3pt}
\SetCell{c,m}\hfil\logo\hfil & \textbf{POLITEKNIK NEGERI MALANG}\par\textbf{JURUSAN TEKNOLOGI INFORMASI}\par\textbf{PROGRAM STUDI: D4 TEKNIK INFORMATIKA} \\
\SetCell[c=2]{c} \textbf{RENCANA PEMBELAJARAN SEMESTER (RPS)} & \\
\end{tblr}
\begin{longtblr}{width=\textwidth,colspec={|Q[l,m,1.9cm]|X[0.85,l,m]|X[1.45,l,m]|X[0.75,l,m]|X[1.15,l,m]|},hlines,vlines,hspan=minimal,rowsep=1pt,colsep=2pt,row{1,3}={bg=headgray,font=\bfseries}}
MATA KULIAH & KODE & BOBOT (SKS)/jam & SEMESTER & TGL. PENYUSUNAN \\
Rekayasa Sistem & RTI256206 & 3 SKS/6 jam & 6 & 1 Juli 2025 \\
OTORISASI & Dosen Pengembang RPS & Koordinator RMK & \SetCell[c=2]{l} Ka PRODI &  \\
 & Agung Nugroho Pramudhita, S.T., M.T. &  & \SetCell[c=2]{l}{\vspace{8mm}\par Dr. Ely Setyo Astuti, S.T., M.T.} &  \\
\SetCell[r=11]{l} Capaian Pembelajaran (CP) & \SetCell[c=4]{l,bg=headgray,font=\bfseries} Capaian Pembelajaran Lulusan Program Studi (CPL-Prodi) &  &  &  \\
 & CPL02 & \SetCell[c=3]{l} Mampu menerapkan prinsip rekayasa sistem dan perangkat lunak untuk menghasilkan solusi aplikasi multi-platform sesuai kebutuhan. &  &  \\
 & CPL06 & \SetCell[c=3]{l} Mampu merancang, mengimplementasikan, menguji, melakukan deployment, memelihara, serta menjamin mutu perangkat lunak yang menjawab permasalahan dan memenuhi kebutuhan pengguna. &  &  \\
 & \SetCell[c=4]{l,bg=headgray,font=\bfseries} Capaian Pembelajaran mata kuliah (CPL-MK) &  &  &  \\
 & CPMK0211 & \SetCell[c=3]{l} Mahasiswa mampu merancang arsitektur sistem skala enterprise yang memenuhi kebutuhan fungsional dan non-fungsional. &  &  \\
 & CPMK0602 & \SetCell[c=3]{l} Mahasiswa mampu melakukan rekayasa kebutuhan (requirement engineering) yang komprehensif untuk sistem kompleks. &  &  \\
 & CPMK0609 & \SetCell[c=3]{l} Mahasiswa mampu merancang dan mengimplementasikan strategi deployment dan operasi sistem yang handal. &  &  \\
 & \SetCell[c=4]{l,bg=headgray,font=\bfseries} Kemampuan akhir tiap tahapan belajar (Sub-CPMK) &  &  &  \\
 & SCPMK0211-05501 & \SetCell[c=3]{l} Mahasiswa mampu merancang arsitektur sistem enterprise meliputi komponen, antarmuka, aliran data, dan ketergantungan antar modul. &  &  \\
 & SCPMK0602-05502 & \SetCell[c=3]{l} Mahasiswa mampu melakukan elicitation, analisis, spesifikasi, dan validasi kebutuhan sistem yang kompleks. &  &  \\
 & SCPMK0609-05503 & \SetCell[c=3]{l} Mahasiswa mampu merancang dan mengimplementasikan pipeline CI/CD serta strategi deployment blue-green/canary untuk sistem produksi. &  &  \\
\SetCell[r=2]{l} Penilaian & \SetCell[c=4]{l} *Nilai CPMK didapat dari agregasi nilai SUB-CPMK yang dibebankan &  &  &  \\
 & \SetCell[c=4]{l}{\tiny\begin{tblr}{width=\linewidth,colspec={|Q[c,m,0.1\linewidth]|Q[c,m,0.16\linewidth]|X[l,m]|X[l,m]|X[l,m]|X[l,m]|X[l,m]|X[l,m]|Q[c,m,0.08\linewidth]|},hlines,vlines,hspan=minimal,rowsep=1pt,colsep=0.7pt,row{1,2}={bg=headgray,font=\bfseries}}
Kode CPMK & Kode SCPMK & \SetCell[c=6]{c} Bobot per Bentuk Penilaian & & & & & & Total Bobot \\
 & & Dok Arsitektur & Sprint 1-2 & UTS & Sprint 3-4 & Laporan SRS & UAS Defense & \\
CPMK0211 & SCPMK0211-05501 & 8\%: desain arsitektur awal & 10\%: implementasi komponen & 0\% & 8\%: finalisasi arsitektur & 0\% & 8\%: defense arsitektur & 34\% \\
CPMK0602 & SCPMK0602-05502 & 7\%: kontribusi SRS awal & 0\% & 15\%: presentasi SRS & 0\% & 10\%: SRS final & 0\% & 32\% \\
CPMK0609 & SCPMK0609-05503 & 0\% & 10\%: pipeline CI/CD & 0\% & 12\%: deployment strategy & 0\% & 12\%: demo deployment & 34\% \\
\SetCell[c=8]{r}\textbf{Total Per Penilaian} & & & & & & & & \textbf{100\%} \\
\end{tblr}} &  &  &  \\
Diskripsi Singkat Mata Kuliah & \SetCell[c=4]{l} Rekayasa Sistem adalah capstone jalur magang/industri yang mengintegrasikan rekayasa kebutuhan, perancangan arsitektur enterprise, dan rekayasa deployment untuk menghasilkan sistem yang handal, maintainable, dan siap produksi. &  &  &  \\
Bahan Kajian / Materi Pembelajaran & \SetCell[c=4]{l}\begin{enumerate}\item Requirement elicitation: teknik wawancara, use case, user story, dan spesifikasi kebutuhan\item Arsitektur sistem enterprise: pola arsitektur, komponen, antarmuka, dan dokumentasi ADR\item Pipeline CI/CD: automated testing, integrasi, dan strategi deployment (blue-green, canary)\item Monitoring, observability, dan pemeliharaan sistem produksi\end{enumerate} &  &  &  \\
\SetCell[r=2]{l} Pustaka & \SetCell[c=4]{l} \textbf{Utama:}\begin{enumerate}\item Sommerville, I. 2016. Software Engineering (10th ed.). Pearson.\item Bass, L. 2013. Software Architecture in Practice. Addison-Wesley.\item Forsgren, N. 2018. Accelerate. IT Revolution.\end{enumerate} &  &  &  \\
 & \SetCell[c=4]{l} \textbf{Pendukung:} Materi LMS, dokumentasi resmi perangkat lunak, dan jobsheet praktikum. &  &  &  \\
Media Pembelajaran & \SetCell[c=2]{l} \textbf{Software:} IDE, Git, CI/CD tools (Jenkins/GitHub Actions), diagramming tools, LMS. &  & \SetCell[c=2]{l} \textbf{Hardware:} Komputer/laptop, proyektor. &  \\
Nama Dosen Pengampu & \SetCell[c=4]{l} Tim Pengajar Program Studi D4 TI &  &  &  \\
Matakuliah Syarat & \SetCell[c=4]{l} - &  &  &  \\
\end{longtblr}
\begin{longtblr}{width=\textwidth,colspec={|Q[c,m,0.06\textwidth]|X[1.35,l,m]|X[1.08,l,m]|X[1.16,l,m]|X[0.7,l,m]|X[1.24,l,m]|X[1.06,l,m]|X[0.98,l,m]|Q[c,m,0.05\textwidth]|},hlines,vlines,hspan=minimal,rowhead=2,row{1,2}={bg=headgreen,font=\bfseries},rowsep=1pt,colsep=1.5pt}
Mgg.\par Ke & Kemampuan Akhir Yang Direncanakan (Sub-CP-MK) & Bahan kajian (Materi Pembelajaran) & Bentuk dan Metode Pembelajaran & Estimasi Waktu & Pengalaman Belajar Mahasiswa & Kriteria \& Bentuk Penilaian & Indikator Penilaian & Bobot Penilaian (\%) \\
(1) & (2) & (3) & (4) & (5) & (6) & (7) & (8) & (9) \\
1 & SCPMK0211-05501 & Pengantar rekayasa sistem: ruang lingkup dan metodologi. & Modalitas: Blended Learning. Bentuk: Luring/Daring. Metode: Case Method, workshop kolaboratif, presentasi teknis, dan peer review. & 1 x 6 x 50' workshop industri; 1 x 6 x 50' dokumentasi dan tugas mandiri. & Mahasiswa mengerjakan aktivitas untuk pengantar rekayasa sistem dan menunjukkan evidence hasil belajar. & Bentuk penilaian: Pengantar Rekayasa Sistem. Mengacu rubrik RTM. & Ketepatan konsep; kualitas artefak; validasi dan dokumentasi. & 3 \\
2 & SCPMK0602-05502 & Requirement elicitation: teknik wawancara dan workshop. & Modalitas: Blended Learning. Bentuk: Luring/Daring. Metode: Case Method, workshop kolaboratif, presentasi teknis, dan peer review. & 1 x 6 x 50' workshop industri; 1 x 6 x 50' dokumentasi dan tugas mandiri. & Mahasiswa mengerjakan aktivitas untuk requirement elicitation dan menunjukkan evidence hasil belajar. & Bentuk penilaian: Requirement Elicitation. Mengacu rubrik RTM. & Ketepatan konsep; kualitas artefak; validasi dan dokumentasi. & 3 \\
3 & SCPMK0602-05502 & Use case, user story, dan spesifikasi kebutuhan. & Modalitas: Blended Learning. Bentuk: Luring/Daring. Metode: Case Method, workshop kolaboratif, presentasi teknis, dan peer review. & 1 x 6 x 50' workshop industri; 1 x 6 x 50' dokumentasi dan tugas mandiri. & Mahasiswa mengerjakan aktivitas untuk use case dan spesifikasi kebutuhan serta menunjukkan evidence hasil belajar. & Bentuk penilaian: Use Case dan Spesifikasi Kebutuhan. Mengacu rubrik RTM. & Ketepatan konsep; kualitas artefak; validasi dan dokumentasi. & 4 \\
4 & SCPMK0211-05501 & Arsitektur sistem: pola dan gaya arsitektur enterprise. & Modalitas: Blended Learning. Bentuk: Luring/Daring. Metode: Case Method, workshop kolaboratif, presentasi teknis, dan peer review. & 1 x 6 x 50' workshop industri; 1 x 6 x 50' dokumentasi dan tugas mandiri. & Mahasiswa mengerjakan aktivitas untuk arsitektur sistem enterprise dan menunjukkan evidence hasil belajar. & Bentuk penilaian: Arsitektur Sistem Enterprise. Mengacu rubrik RTM. & Ketepatan konsep; kualitas artefak; validasi dan dokumentasi. & 4 \\
5 & SCPMK0211-05501 & Sprint 1: desain komponen dan antarmuka sistem. & Modalitas: Blended Learning. Bentuk: Luring/Daring. Metode: Case Method, workshop kolaboratif, presentasi teknis, dan peer review. & 1 x 6 x 50' workshop industri; 1 x 6 x 50' dokumentasi dan tugas mandiri. & Mahasiswa mengerjakan aktivitas sprint 1 untuk desain komponen dan antarmuka sistem serta menunjukkan evidence hasil belajar. & Bentuk penilaian: Sprint 1 Desain Komponen. Mengacu rubrik RTM. & Ketepatan konsep; kualitas artefak; validasi dan dokumentasi. & 5 \\
6 & SCPMK0211-05501 & Sprint 1: prototipe dan validasi arsitektur. & Modalitas: Blended Learning. Bentuk: Luring/Daring. Metode: Case Method, workshop kolaboratif, presentasi teknis, dan peer review. & 1 x 6 x 50' workshop industri; 1 x 6 x 50' dokumentasi dan tugas mandiri. & Mahasiswa mengerjakan aktivitas sprint 1 untuk prototipe dan validasi arsitektur serta menunjukkan evidence hasil belajar. & Bentuk penilaian: Sprint 1 Prototipe dan Validasi Arsitektur. Mengacu rubrik RTM. & Ketepatan konsep; kualitas artefak; validasi dan dokumentasi. & 5 \\
7 & SCPMK0602-05502 & Sprint 1 review dan SRS validation meeting. & Modalitas: Blended Learning. Bentuk: Luring/Daring. Metode: Case Method, workshop kolaboratif, presentasi teknis, dan peer review. & 1 x 6 x 50' workshop industri; 1 x 6 x 50' dokumentasi dan tugas mandiri. & Mahasiswa mengerjakan aktivitas sprint 1 review dan SRS validation meeting serta menunjukkan evidence hasil belajar. & Bentuk penilaian: Sprint 1 Review dan SRS Validation. Mengacu rubrik RTM. & Ketepatan konsep; kualitas artefak; validasi dan dokumentasi. & 5 \\
8 & SCPMK0602-05502 & UTS Milestone: presentasi arsitektur dan SRS. & Modalitas: Blended Learning. Bentuk: Luring/Daring. Metode: Case Method, workshop kolaboratif, presentasi teknis, dan peer review. & 1 x 6 x 50' workshop industri; 1 x 6 x 50' dokumentasi dan tugas mandiri. & Mahasiswa mengerjakan aktivitas UTS berupa presentasi arsitektur dan SRS serta menunjukkan evidence hasil belajar. & Bentuk penilaian: UTS Presentasi Arsitektur dan SRS. Mengacu rubrik RTM. & Ketepatan konsep; kualitas artefak; validasi dan dokumentasi. & 15 \\
9 & SCPMK0609-05503 & Sprint 2: implementasi pipeline CI/CD. & Modalitas: Blended Learning. Bentuk: Luring/Daring. Metode: Case Method, workshop kolaboratif, presentasi teknis, dan peer review. & 1 x 6 x 50' workshop industri; 1 x 6 x 50' dokumentasi dan tugas mandiri. & Mahasiswa mengerjakan aktivitas sprint 2 untuk implementasi pipeline CI/CD dan menunjukkan evidence hasil belajar. & Bentuk penilaian: Sprint 2 Pipeline CI/CD. Mengacu rubrik RTM. & Ketepatan konsep; kualitas artefak; validasi dan dokumentasi. & 4 \\
10 & SCPMK0609-05503 & Sprint 2: automated testing dan integrasi. & Modalitas: Blended Learning. Bentuk: Luring/Daring. Metode: Case Method, workshop kolaboratif, presentasi teknis, dan peer review. & 1 x 6 x 50' workshop industri; 1 x 6 x 50' dokumentasi dan tugas mandiri. & Mahasiswa mengerjakan aktivitas sprint 2 untuk automated testing dan integrasi serta menunjukkan evidence hasil belajar. & Bentuk penilaian: Sprint 2 Automated Testing dan Integrasi. Mengacu rubrik RTM. & Ketepatan konsep; kualitas artefak; validasi dan dokumentasi. & 5 \\
11 & SCPMK0609-05503 & Sprint 2 review dan deployment strategy planning. & Modalitas: Blended Learning. Bentuk: Luring/Daring. Metode: Case Method, workshop kolaboratif, presentasi teknis, dan peer review. & 1 x 6 x 50' workshop industri; 1 x 6 x 50' dokumentasi dan tugas mandiri. & Mahasiswa mengerjakan aktivitas sprint 2 review dan deployment strategy planning serta menunjukkan evidence hasil belajar. & Bentuk penilaian: Sprint 2 Review dan Deployment Strategy. Mengacu rubrik RTM. & Ketepatan konsep; kualitas artefak; validasi dan dokumentasi. & 4 \\
12 & SCPMK0609-05503 & Sprint 3: konfigurasi staging dan production environment. & Modalitas: Blended Learning. Bentuk: Luring/Daring. Metode: Case Method, workshop kolaboratif, presentasi teknis, dan peer review. & 1 x 6 x 50' workshop industri; 1 x 6 x 50' dokumentasi dan tugas mandiri. & Mahasiswa mengerjakan aktivitas sprint 3 untuk konfigurasi staging dan production environment serta menunjukkan evidence hasil belajar. & Bentuk penilaian: Sprint 3 Staging dan Production Environment. Mengacu rubrik RTM. & Ketepatan konsep; kualitas artefak; validasi dan dokumentasi. & 5 \\
13 & SCPMK0609-05503 & Sprint 3: blue-green deployment dan rollback strategy. & Modalitas: Blended Learning. Bentuk: Luring/Daring. Metode: Case Method, workshop kolaboratif, presentasi teknis, dan peer review. & 1 x 6 x 50' workshop industri; 1 x 6 x 50' dokumentasi dan tugas mandiri. & Mahasiswa mengerjakan aktivitas sprint 3 untuk blue-green deployment dan rollback strategy serta menunjukkan evidence hasil belajar. & Bentuk penilaian: Sprint 3 Blue-Green Deployment. Mengacu rubrik RTM. & Ketepatan konsep; kualitas artefak; validasi dan dokumentasi. & 4 \\
14 & SCPMK0609-05503 & Monitoring, logging, dan observability. & Modalitas: Blended Learning. Bentuk: Luring/Daring. Metode: Case Method, workshop kolaboratif, presentasi teknis, dan peer review. & 1 x 6 x 50' workshop industri; 1 x 6 x 50' dokumentasi dan tugas mandiri. & Mahasiswa mengerjakan aktivitas untuk monitoring, logging, dan observability sistem serta menunjukkan evidence hasil belajar. & Bentuk penilaian: Monitoring Logging Observability. Mengacu rubrik RTM. & Ketepatan konsep; kualitas artefak; validasi dan dokumentasi. & 5 \\
15 & SCPMK0211-05501 & Dokumentasi sistem: ADR dan arsitektur diagram. & Modalitas: Blended Learning. Bentuk: Luring/Daring. Metode: Case Method, workshop kolaboratif, presentasi teknis, dan peer review. & 1 x 6 x 50' workshop industri; 1 x 6 x 50' dokumentasi dan tugas mandiri. & Mahasiswa mengerjakan aktivitas untuk dokumentasi sistem ADR dan arsitektur diagram serta menunjukkan evidence hasil belajar. & Bentuk penilaian: Dokumentasi ADR dan Arsitektur Diagram. Mengacu rubrik RTM. & Ketepatan konsep; kualitas artefak; validasi dan dokumentasi. & 5 \\
16 & SCPMK0602-05502 & SRS final dan sign-off dengan stakeholder. & Modalitas: Blended Learning. Bentuk: Luring/Daring. Metode: Case Method, workshop kolaboratif, presentasi teknis, dan peer review. & 1 x 6 x 50' workshop industri; 1 x 6 x 50' dokumentasi dan tugas mandiri. & Mahasiswa mengerjakan aktivitas untuk SRS final dan sign-off dengan stakeholder serta menunjukkan evidence hasil belajar. & Bentuk penilaian: SRS Final dan Sign-Off. Mengacu rubrik RTM. & Ketepatan konsep; kualitas artefak; validasi dan dokumentasi. & 4 \\
17 & SCPMK0211-05501, SCPMK0609-05503 & UAS: defense arsitektur dan demo deployment pipeline. & Modalitas: Blended Learning. Bentuk: Luring/Daring. Metode: Case Method, workshop kolaboratif, presentasi teknis, dan peer review. & 1 x 6 x 50' workshop industri; 1 x 6 x 50' dokumentasi dan tugas mandiri. & Mahasiswa mengerjakan aktivitas UAS berupa defense arsitektur dan demo deployment pipeline serta menunjukkan evidence hasil belajar. & Bentuk penilaian: UAS Defense Arsitektur dan Demo Deployment. Mengacu rubrik RTM. & Ketepatan konsep; kualitas artefak; validasi dan dokumentasi. & 20 \\
\end{longtblr}
